\section{Strict Gaussian contraction and the diffuse case}\label{sec:gauss}
Let $\gamma_d$ be standard Gaussian measure and
$Q_rv(x)=\E v(rx+\sqrt{1-r^2}Z)$. Define $\cI_\gamma$ by the same
entropy formula as on the cube.

\begin{lemma}\label{lem:gaussstrict}
For every measurable nonconstant $v:\R\to[-1,1]$ with
$0<\cI_\gamma(v)<L$ and every $0<r<1$,
\[
 \cI_\gamma(Q_rv)<f_r(\cI_\gamma(v)).
\]
\end{lemma}
\begin{proof}
First the non-strict inequality transfers from the assumed cube theorem.
For smooth bounded $v$ with bounded third derivative apply it to the
normalized sum $S_N=N^{-1/2}\sum X_i$. Conditional on the whole noisy
output $Y=y$, the centered input summands have variance $(1-r^2)/N$
and third absolute moments $O(N^{-3/2})$. Successive Gaussian replacement,
with second-order Taylor cancellation, gives
\[
 \sup_y|T_r(v\circ S_N)(y)-Q_rv(S_N(y))|
 \le C_r\|v'''\|_\infty N^{-1/2}.
\]
The usual central limit theorem and uniform continuity of $\psi$ now
transfer both input and output information. Gaussian smoothing approximates
any bounded measurable kernel in $L^2$, and $Q_r$ is a contraction,
so the non-strict bound follows for every such kernel. Independent blocks
give the same statement in each fixed Gaussian dimension.

We need the following scalar inequality. If $0<a,r<1$,
$k=rV(ra)/V(a)$, and $b=e^{\psi(a)}-1$, then
\begin{equation}\label{eq:scalarstrict}
 0<k<r^2<1,\qquad k>\frac{r^2(1-b)}{1-r^2b}.
\end{equation}
Here are details of the strict comparison. Put
$M_0(a)=[a/(1-a^2)-V(a)]/[a/(1-a^2)+V(a)]$ and
\[
 K_0(a)=\psi(a)+\log[a+(1-a^2)V(a)]-\log(2a).
\]
The latter has zero limits at both endpoints. Its derivative has the sign
of $1-t\coth t+t^2\operatorname{sech}^2t$, where $a=\tanh t$.
The partial fraction for $\coth$ \cite[4.36.3]{DLMF} gives
\[
 \frac{t\coth t-1}{t^2\operatorname{sech}^2t}
 =2\sum_{j\ge1}\frac{\cosh^2t}{t^2+\pi^2j^2}.
\]
Each summand strictly increases, since
$\tanh t>t/(1+t^2)>t/(t^2+\pi^2j^2)$; the first inequality follows
by differentiating $(1+t^2)\sinh t-t\cosh t$.
The sum increases from $1/3$ to infinity. Thus $K_0$ first increases
and then decreases, so $K_0(a)>0$, equivalently $b>M_0(a)$.
Moreover $M_0(z)/z^2$ increases: its numerator and denominator series
have coefficients $(2j+2)/(2j+3)$ and $(2j+2)/(2j+1)$, whose ratios
strictly increase. Pairing the quotient-rule numerator over $j>i$
proves that sign directly. If $A(x)=V(a\sqrt x)/(a\sqrt x)$, then
\[
 \frac{A'(x)}{A(x)+xA'(x)}=\frac{M_0(a\sqrt x)}x\le M_0(a)<b.
\]
Hence $(1-bx)A(x)$ decreases; comparison at $r^2$ and one proves
\eqref{eq:scalarstrict}. Also $f_r'>0$ and $f_r''<0$, because
$aV'(a)/V(a)=\sinh(2t)/(2t)$ strictly increases.

Suppose first that $w=Q_\tau v$, $0<\tau<1$, is nonconstant and
attains equality. Write $h=Q_rw$, $m=\E w$, $\cI_\gamma(w)=\psi(a)$.
The likelihood ratio of a binary output of conditional mean $w$ relative
to its product marginals is $(1+sw)/(1+sm)$, $s=\pm1$.
Its product second moment is $1+q$, where
$q=\Var(w)/(1-m^2)$. Jensen under the joint measure gives
\begin{equation}\label{eq:likelihood}
 \cI_\gamma(w)\le\log(1+q),\qquad q\ge b.
\end{equation}

We record the regularity needed for the variation. Both $1\pm v$ have
positive integrals on some finite interval. Their Gaussian convolutions
therefore satisfy $1\pm w(x)\ge c\exp(-Cx^2)$; hence
$|V(w(x))|\le C(1+x^2)$. All ordinary derivatives of $w$ are bounded
by differentiating the convolution kernel. With $A_s=1+sw$ and
$g_s=1+sv$, its score formula is
$A_s'=(\tau/\sqrt{1-\tau^2})\E[Zg_s(\tau x+\sqrt{1-\tau^2}Z)]$.
Weighted Cauchy--Schwarz, averaged over $s=\pm1$, gives
\begin{equation}\label{eq:fisher}
 \frac{w'^2}{1-w^2}\le\frac{\tau^2}{1-\tau^2}.
\end{equation}
For the tilted density $g_s/A_s$ its entropy is at most
$\log(2/A_s)$. The Gaussian entropy variational inequality
$\lambda\E_{\rm tilt}Z-\lambda^2/2\le\operatorname{Ent}(g_s/A_s)$,
which follows by comparison with the exponential density, bounds the
tilted mean by $\sqrt{2\log(2/A_s)}$.
Consequently $|(V\circ w)'(x)|\le C(1+|x|)$.
The same bounds hold for $h$.

The nonnegative functional
$\mathcal G(v)=f_r(\cI_\gamma(v))-\cI_\gamma(Q_rv)$ vanishes at $w$.
A smooth compactly supported perturbation $\phi$ obeys
$|\phi|\le C(1\pm w)$; positivity of $Q_r$ gives
$|Q_r\phi|\le C(1\pm h)$. Thus perturbations of both signs are
admissible and dominated differentiation of entropy is justified by the
quadratic growth above. The first variation and self-adjointness give
\begin{equation}\label{eq:euler}
 Q_rV(h)=kV(w)+(1-k)V(m).
\end{equation}
The terms are continuous, so this identity is pointwise.
Differentiation, the commutation $(Q_rf)'=rQ_rf'$, and multiplication
by $w'$ give
\begin{equation}\label{eq:energygauss}
 r^2\E\frac{(Q_rw')^2}{1-h^2}=k\mathcal D,
 \qquad \mathcal D=\E\frac{w'^2}{1-w^2}>0.
\end{equation}
All integrals and the self-adjoint interchange are legitimate: the
composed derivative of $V(h)$ has linear growth, $w'$ is bounded,
and \eqref{eq:fisher} controls the Fisher integrals.

Translate $w_c(x)=w(x+c)$. Its noisy output is $h(x+c/r)$.
The same bounds, locally uniformly in $c$, allow two entropy derivatives.
Put $A=\E[(V(w)-V(m))w']$. The linear terms involving $w''$
cancel by \eqref{eq:euler}, leaving
\[
 \frac{d^2}{dc^2}\mathcal G(w_c)\bigg|_{c=0}
 =k\mathcal D+f_r''(\psi(a))A^2
 -\E\frac{(Q_rw')^2}{1-h^2}
 +(1-k)\frac{(\E w')^2}{1-m^2}.
\]
Since $(\E w')^2\le\mathcal D(1-m^2)(1-q)$,
\eqref{eq:energygauss}, \eqref{eq:likelihood}, and strict concavity give
\[
 \mathcal G''(0)\le\mathcal D
 \{k(1-r^{-2})+(1-k)(1-b)\}<0
\]
by \eqref{eq:scalarstrict}. This contradicts minimality.
For an arbitrary equality kernel choose $r<\tau<1$, $s=r/\tau$.
The non-strict inequality and the identity $f_s\circ f_\tau=f_r$ force
$Q_\tau v$ to attain equality at $s$, with positive information below
$L$. This is impossible by the smoothed case.
\end{proof}

\begin{corollary}\label{cor:gaussgap}
For every compact rectangle $0<b_0\le b\le b_1<1$,
$0<\sigma_0\le\sigma\le\sigma_1<1$, the deficit
$\psi(\sigma b)-\cI_\gamma(Q_\sigma v)$ has a strictly positive
lower bound over all finite Gaussian dimensions and bounded kernels
with $\cI_\gamma(v)=\psi(b)$.
\end{corollary}
\begin{proof}
Put $f=(1+v)/2$ and let $f^*$ be its increasing scalar Gaussian quantile.
For an output set $A$ of measure $p$, layer cake and Borell's two-set
inequality \cite[Theorem 1.1]{MN} give
\[
 \int_A Q_\sigma f\,d\gamma
 \le\int_{A^*}Q_\sigma f^*\,d\gamma,
\]
where $A^*$ is the upper ray of measure $p$. The level sets of the same
$f^*$ are nested rays, so all comparisons are simultaneous.
Taking largest-set integrals and then the supremum over $p$ of their
value minus $cp$ gives the stop-loss inequalities for every $c$.
Equal means and convex piecewise-linear approximation give convex order.
Apply it to $t\mapsto\psi(2t-1)$: the rearrangement preserves input
information and does not decrease output information.
The class of bounded nondecreasing scalar kernels is compact in
$L^2(\gamma)$ by Helly selection and dominated convergence.
Information is continuous there. Joint continuity of $Q_\sigma$ in
kernel and correlation follows from contraction and the Hermite
multipliers for the fixed limit kernel. The constrained rectangle is
compact, and its continuous deficit is positive by
Lemma~\ref{lem:gaussstrict}. Its minimum proves the claim.
\end{proof}

\subsection{Transfer from increasing cube kernels}
Suppose increasing kernels $g_j$, possibly in changing dimensions,
satisfy $\Delta_{r_j}(g_j)\to0$ while all singletons are bounded by
$\alpha_j\to0$. Fix $r_1<\lambda<1$, put $h_j=T_\lambda g_j$,
and $\sigma_j=r_j/\lambda$. Equation~\eqref{eq:budget} gives
\[
 \cI(h_j)=\psi(\lambda a_j)+o(1),\qquad
 \cI(T_{\sigma_j}h_j)=\psi(r_ja_j)+o(1).
\]
For a fixed degree $d$, let $p_j=\Pi_{\le d}h_j$ and use the same
multilinear polynomial on Gaussian variables. Its cube tail has norm
at most $\lambda^{d+1}$.
Monotonicity gives $0\le\partial_i h_j\le1$, so
\[
 \operatorname{Inf}_i(p_j)\le\E(\partial_i h_j)^2
 \le\E\partial_i h_j=\lambda\widehat g_j(\{i\})\le\alpha_j,
 \qquad \sum_i\operatorname{Inf}_i(p_j)\le d.
\]
For a $C^4$ test function with bounded fourth derivative, replacing one
sign at a time by a Gaussian cancels its first three moments. The
fourth-order error sums to
$C_d\|\varphi^{(4)}\|_\infty\sum_i\operatorname{Inf}_i(p_j)^2
\le C_d d\alpha_j\|\varphi^{(4)}\|_\infty$.
Here the needed mixed-product fourth moment is
$\E B^4\le9^{d-1}(\E B^2)^2$ for the derivative polynomial $B$;
it follows by tensorizing the two-point/Gaussian inequality
$\E(A+XB/\sqrt3)^4\le(A^2+B^2)^2$.
This is the fixed-degree replacement argument of \cite{MOO}, applied
also to $q_j=Q_{\sigma_j}p_j$, whose coefficients are multiplied by
$\sigma_j^{|S|}$.

The bounded tests $\clip(t)$ and $\psi(\clip(t))$, where clipping is
to $[-1,1]$, are handled by uniform smooth approximation.
For the unbounded test $D(t)=\dist(t,[-1,1])^2$, mollify at radius
$\tau$. Its fourth derivative is bounded for each $\tau>0$, while
$|D_\tau(t)-D(t)|\le C\tau(1+|t|)+C\tau^2$.
The second moments of the polynomials equal their coefficient norms
and are at most one under both measures. Thus first $j\to\infty$,
then $\tau\downarrow0$, transfers this test as well.
Under cube measure $\E D(p_j)\le\lambda^{2d+2}$, hence the same
Gaussian expectation is at most $\lambda^{2d+2}+o(1)$.
Set $v_j=\clip(p_j)$ in Gaussian space. Since $Q_{\sigma_j}v_j$ is
bounded, contraction and the metric projection property give
\[
 \|v_j-p_j\|_2\le\lambda^{d+1}+o(1),\qquad
 \|Q_{\sigma_j}v_j-\clip(q_j)\|_2\le\|v_j-p_j\|_2.
\]
Together with the bounded test comparisons and uniform entropy
continuity, these transfer both information budgets with an error
that tends to zero as $d\to\infty$, uniformly after $j\to\infty$.
Choose the Gaussian gap in Corollary~\ref{cor:gaussgap} on a slightly
larger compact rectangle around $\lambda[a_0,a_1]$ and
$[r_0/\lambda,r_1/\lambda]$ first. Then choose $d$ so all tail errors
are smaller than a fixed fraction of that gap, and finally choose $j$.
At the exact amplitude $b_j=\psi^{-1}(\cI_\gamma(v_j))$, the transferred
deficit contradicts that gap.

Therefore there are $\beta,\varepsilon>0$ such that every increasing
kernel with deficit below $\varepsilon$ has a singleton at least $\beta$.
Lemma~\ref{lem:marked} handles those kernels; above the threshold use
$d_a\le4$. Sorting and \eqref{eq:compression} prove
Theorem~\ref{thm:compact}.

\subsection{Sharpness of the deficit power}
Fix $a,r\in(0,1)$ and use the two-dimensional kernel
$u_\epsilon=c(\epsilon)X_1+\epsilon X_2$, normalized by
$\cI(u_\epsilon)=\psi(a)$. The implicit function theorem gives
\[
 c(\epsilon)=a-\frac{\epsilon^2}{2(1-a^2)V(a)}+O(\epsilon^4).
\]
It is interior for small $\epsilon$, has distance
$d_a=\epsilon^2+O(\epsilon^4)$, and
\[
 \Delta_r(u_\epsilon)=\frac12
 \left\{\frac{rV(ra)}{(1-a^2)V(a)}
             -\frac{r^2}{1-r^2a^2}\right\}\epsilon^2+O(\epsilon^4).
\]
The coefficient is positive by Lemma~\ref{lem:local}.
Thus no larger power of the deficit can give a uniform stability bound.
